% 自動生成（tools/mknext.py p1）。各章の道しるべ「この章の先に」から、後ろの巻の章を集めたもの（第13回）
\begin{longtblr}{colspec={Q[l,m,wd=58mm] X[l,m]}, width=\linewidth, rowhead=1, row{1}={bg=bkVol!14, font=\sffamily\bfseries\footnotesize},
  hlines={0.4pt,bkLine}, rowsep=1.5pt, cells={font=\footnotesize}}
この巻の章 & この道具を使う章 \\
{\bknextpill{力-2}\ 力の図から式へ：運動方程式を立てる} & {\bknextpill{波-1}\ ばねとおもりで考える：1自由度の振動} \\
{\bknextpill{力-3}\ ぶつかる・噴き出す：運動量と衝突} & {\bknextpill{熱-4}\ 流れる気体の熱：エンタルピーとノズル} \\
{\bknextpill{力-4}\ 高さが速さに変わる：仕事とエネルギー} & {\bknextpill{熱-2}\ 仕事で温まる：第一法則と内部エネルギー} \\
{\bknextpill{力-5}\ 橋を支える三角形：つり合いとトラス} & {\bknextpill{電-5}\ 回路を行列で解く：直流回路の解析} \\
{\bknextpill{力-6}\ 回る物体：慣性モーメントと角運動量} & {\bknextpill{電-12}\ モーターと変圧器のしくみ：電気機器} \\
{\bknextpill{力-7}\ 転がる・ロボットの腕：平面運動と機構} & {\bknextpill{波-3}\ つながるゆれ：多自由度振動とラグランジュの式\\ \bknextpill{波-4}\ ねらいどおりに動かす：フィードバック制御} \\
{\bknextpill{力-8}\ 引っぱる・押す：応力とひずみ} & {\bknextpill{熱-1}\ 熱の出入りを数える：温度・比熱・潜熱\\ \bknextpill{波-5}\ 波を式で書く：進行波と波動方程式} \\
{\bknextpill{力-9}\ 向きで変わる応力：組み合わせ応力とモールの円} & {\bknextpill{波-10}\ 見えない中を調べる：超音波と振動の計測\\ \bknextpill{電-4}\ 電気の流れと熱：電流と抵抗} \\
{\bknextpill{力-10}\ 板がたわむ：はりの曲げとたわみ} & {\bknextpill{波-6}\ はね返って重なる：反射・定在波・固有振動} \\
{\bknextpill{力-11}\ ねじる・ふくらむ・曲がる：ねじりと座屈} & {\bknextpill{熱-3}\ ふくらむ・縮む気体：理想気体の状態変化\\ \bknextpill{波-1}\ ばねとおもりで考える：1自由度の振動} \\
{\bknextpill{力-12}\ 水の中の力：圧力と浮力} & {\bknextpill{熱-3}\ ふくらむ・縮む気体：理想気体の状態変化} \\
{\bknextpill{力-13}\ 流れを数える：質量・エネルギー・運動量} & {\bknextpill{熱-4}\ 流れる気体の熱：エンタルピーとノズル} \\
{\bknextpill{力-14}\ 管の中の流れ：粘性・層流と乱流・圧力損失} & {\bknextpill{熱-12}\ 流れが熱を運ぶ：対流熱伝達と沸騰} \\
{\bknextpill{力-15}\ 模型で本物を調べる：物体まわりの流れと相似則} & {\bknextpill{熱-4}\ 流れる気体の熱：エンタルピーとノズル\\ \bknextpill{熱-12}\ 流れが熱を運ぶ：対流熱伝達と沸騰} \\
{\bknextpill{力-16}\ コンピュータで解く力学} & {\bknextpill{熱-11}\ 壁を通りぬける熱：熱伝導と熱抵抗\\ \bknextpill{波-3}\ つながるゆれ：多自由度振動とラグランジュの式\\ \bknextpill{波-5}\ 波を式で書く：進行波と波動方程式\\ \bknextpill{電-5}\ 回路を行列で解く：直流回路の解析} \\
\end{longtblr}
